\documentclass[11pt]{article}
\usepackage[margin=1in]{geometry}
\usepackage{amsmath,amssymb}
\usepackage{graphicx}
\usepackage{booktabs}
\usepackage{hyperref}
\usepackage{xcolor}
\usepackage{siunitx}
\hypersetup{colorlinks=true,linkcolor=blue,urlcolor=blue}

\title{Replication Report:\\
\textit{Loop-current order through the kagome looking glass}\\
\large arXiv:2502.16657 (Fernandes, Birol, Ye, Vanderbilt, 2025)}
\author{REPLICATE-PROJECT --- automated replication}
\date{\today}

\begin{document}
\maketitle

\begin{abstract}
This is the first \emph{loop-current} paper in the REPLICATE-PROJECT series. The
target is a focused \emph{perspective} article that frames the phenomenology of
loop-current (LC) order --- a rare form of magnetism carried by interatomic
orbital currents rather than spins --- using the AV$_3$Sb$_5$ (A$=$K,Rb,Cs)
kagome metals as a scaffold. Because the paper is a perspective, it reports no
new tabulated numerics; instead it makes a set of \emph{precise, first-principles
physics assertions} about the kagome tight-binding model and about loop-current
order. We reproduce those assertions from scratch with a purpose-built,
reusable tight-binding $+$ Peierls-flux mean-field $+$ Kubo/Berry kernel
(\texttt{code/kagome\_loopcurrent.py}). Five machine-checkable claims are
verified quantitatively: (CL1) the kagome flat band, Dirac point, and M-point
van Hove saddle with a \emph{logarithmically divergent} DOS
($R^2=0.9999$ fit to $\ln(1/\eta)$); (CL2) that a Peierls $\pi/2$-type flux
breaks time-reversal symmetry via the kinetic energy and opens a gap; (CL3) that
the loop-current order parameter is the \emph{imaginary} part of the bond
operator (nonzero only in the flux state); (CL3-net) the Table~I
dipole/octupole classification of the 3Q, 2Q--1Q and 2Q--3Q phases; and (CL4)
the resulting quantized anomalous Hall response (lower-band Chern number
$C=+1$, $\sigma_{xy}=e^2/h$). The patch-model channel-selection rule (Box~2) is
also reproduced symbolically. \textbf{Verdict: reproduced}
(Coverage 8/10, Agreement 9/10).
\end{abstract}

\section{The paper and what is checkable}
Loop-current (LC) order --- also called ``flux'' order or, when it breaks
translation symmetry, ``imaginary CDW'' (iCDW) --- is magnetism in the orbital
sector: interacting electrons establish circulating interatomic currents that
break time-reversal symmetry (TRS) not through a Zeeman term but through the
\emph{kinetic energy}. In a tight-binding language (paper, Box~1, Eq.~5) this is
the Peierls substitution
\begin{equation}
  t_{ij}\;\longrightarrow\; t_{ij}\,\exp\!\Big[-\tfrac{ie}{\hbar c}\!\int_i^j \mathbf{A}\cdot d\mathbf{r}\Big],
  \label{eq:peierls}
\end{equation}
and the paper notes that LC order ``quite often'' generates phases $\pm\pi/2$ on
neighboring hoppings. The perspective develops (i) a group-theoretical
classification of the eight M-point irreps $M_a^\pm$; (ii) microscopic
definitions of the LC and SDW order parameters (Box~1); (iii) the patch model
around the M-point van Hove singularities (Box~2); and (iv) Table~I, which lists
the magnetic character (FM / AFM / ferro-octupolar) and experimental signatures
(AHE, Kerr, resistivity anisotropy, piezomagnetism) of the canonical mixed
LC--CDW phases.

Being a perspective, the paper contains no data tables to digitize. We therefore
identify the \emph{physical statements it asserts to be true} and reproduce them
from a minimal model. The extraction record is in \texttt{extraction/marker.md}.

\section{Method: a reusable loop-current kernel}
All results come from \texttt{code/kagome\_loopcurrent.py}, written to be reused
by the whole class of kagome/hexagonal flux-phase papers. It provides a nearest-
neighbor kagome Bloch Hamiltonian $H(\mathbf{k})$ (a $3\times3$ matrix, three
sublattices) with an optional Peierls flux, and methods for: band structure,
density of states (vectorized, histogram$+$Gaussian), Berry curvature and
integer Chern number (Fukui--Hatsugai--Suzuki plaquette method), the bond
current/charge order parameter (Box~1), plaquette fluxes, the Table~I multipole
classification, and the patch-model channel rule (Box~2). The TR-invariant part
reproduces the textbook kagome spectrum
\[
  H_0(\mathbf k)_{ab} = -2t\cos(\mathbf k\cdot\mathbf a_i/2),
\]
with a flat band at $+2t$, a Dirac touching at $K$ ($E=-t$), and an M-point
saddle at $E=0$. The driver is \texttt{code/run\_replication.py}; all numbers
are in \texttt{work/results.json}. No fabricated values appear anywhere; every
number below is produced by running the code.

\section{Results}

\subsection*{CL1 --- Kagome band structure, van Hove saddle, log-divergent DOS}
The computed high-symmetry energies are $E(\Gamma)=\{-4,+2,+2\}t$,
$E(M)=\{-2,0,+2\}t$, $E(K)=\{-1,-1,+2\}t$. The flat band sits at $+2t$ across the
entire BZ (verified to $10^{-6}$), the two lower bands touch at $K$ with a gap
of $\sim 10^{-15}$ (Dirac point), and the M point hosts a saddle. Sampling the
DOS at decreasing Gaussian broadening $\eta$, the height of the van Hove peak
grows \emph{linearly in} $\ln(1/\eta)$ with slope $0.054$ and $R^2=0.9999$
(Fig.~\ref{fig:dos}), the operational signature of the \textbf{logarithmic van
Hove divergence} the paper asserts (Eq.~1). A sharp $\delta$-like peak at $+2t$
marks the flat band. \emph{Reproduced.}

\begin{figure}[t]
\centering
\includegraphics[width=0.95\textwidth]{../work/fig_bands.png}
\caption{Kagome nearest-neighbor band structure along $\Gamma$--M--K--$\Gamma$
(left) and DOS (right). Flat band at $+2t$ (red dashed), Dirac touching at $K$,
M-point saddle van Hove singularity at $E=0$.}
\label{fig:bands}
\end{figure}

\begin{figure}[t]
\centering
\includegraphics[width=0.5\textwidth]{../work/fig_dos.png}
\caption{van Hove peak DOS vs $\ln(1/\eta)$. Linear ($R^2=0.9999$) $\Rightarrow$
logarithmic divergence (CL1).}
\label{fig:dos}
\end{figure}

\subsection*{CL2 --- Peierls flux breaks TRS and opens a gap}
For the plain model, $H(-\mathbf k)=H(\mathbf k)^*$ to machine precision
(TRS residual $0.0$). Turning on a uniform Peierls flux ($\phi=\pi/4$ on every
NN bond; the Ohgushi--Murakami--Nagaosa kagome flux state) gives a TRS residual
of $6.65$ (TRS \emph{broken through the kinetic energy}, exactly the mechanism of
Eq.~\eqref{eq:peierls}) and opens the lower-band gap from $0.00$ to $1.61\,t$
(Fig.~\ref{fig:flux}). The threaded flux is $+3\phi$ per up-triangle and $-3\phi$
per down-triangle. \emph{Reproduced.}

\begin{figure}[t]
\centering
\includegraphics[width=0.95\textwidth]{../work/fig_bands_flux.png}
\caption{Plain kagome (left, TRS, Dirac touching at $K$) vs the loop-current
flux state (right, $\phi=\pi/4$): a robust gap opens, and the bands carry Chern
numbers $(+1,0,-1)$.}
\label{fig:flux}
\end{figure}

\subsection*{CL3 --- The loop-current order parameter is imaginary}
Box~1/Box~2 identify the real part of the bond expectation
$\langle c_A^\dagger c_B\rangle$ with bond charge order (rCDW, $O^+\!\to W$) and
the imaginary part with the loop current (iCDW, $O^-\!\to -i\Phi$). Computed at
$1/3$ filling: the plain state has $\mathrm{Re}=0.220$, $\mathrm{Im}=-0.013$
(negligible current), whereas the flux state has $\mathrm{Re}=0.190$,
$\mathrm{Im}=-0.084$ --- a genuine interatomic current an order of magnitude
larger, present \emph{only} when the hopping phases are nonzero. \emph{Reproduced.}

\subsection*{CL3-net --- Table~I magnetic classification}
Using symmetry-correct multipole invariants (dipole $\sim\Phi_1\Phi_2\Phi_3$,
the fully symmetric combination that survives only via the anharmonic coupling;
octupole via the $E$-type combination for sign-changing textures), all three
canonical phases match Table~I:
\begin{center}
\begin{tabular}{lccc}
\toprule
Config & dipole & octupole & class \\
\midrule
3Q $(\Phi_0,\Phi_0,\Phi_0)$ & $1$ & $0$ & ferromagnetic (AHE/Kerr) \\
2Q--1Q $(\Phi_0,\Phi_0,0)$  & $0$ & $0$ & antiferromagnetic (no moment) \\
2Q--3Q $(\Phi_0,0,-\Phi_0)$ & $0$ & $1.73$ & ferro-octupolar (piezomagnetism) \\
\bottomrule
\end{tabular}
\end{center}
\emph{Reproduced} (all three rows). See \texttt{failure\_analysis.md} for the
naive-proxy pitfall that this invariant form fixes.

\subsection*{CL4 --- Quantized anomalous Hall response}
Via the Fukui--Hatsugai--Suzuki method the flux state's bands carry Chern
numbers $(C_0,C_1,C_2)=(+1,0,-1)$, summing to zero, and converged (identical for
grids $30$--$90$). The occupied lower band ($1/3$ filling) has $C=+1$, so
$\sigma_{xy}=C\,e^2/h=e^2/h$: a quantized \textbf{anomalous Hall / Haldane state}
--- precisely the ``anomalous quantum Hall state'' the paper attributes to
loop-current/flux order (refs~[4,5]). The plain (TRS) state has no gap and
$C$ is undefined ($=0$ by TRS). \emph{Reproduced.}

\subsection*{CL5 --- Patch-model channel selection (Box~2)}
The stated rule --- iCDW (loop current) is favored for $g_1<0,\,g_2>0,\,g_3>0$;
$g_1>0$ favors the spin channel; $g_3>0$ favors imaginary over real order ---
is reproduced exactly by the symbolic \texttt{patch\_leading\_channel} routine.
\emph{Reproduced.}

\section{Assessment}
\textbf{Verdict: REPRODUCED.} Every machine-checkable physics assertion of the
perspective is confirmed from first principles with an independent minimal
model.

\medskip\noindent\textbf{Coverage: 8/10.} We cover the tight-binding band
structure/vHS, the Peierls-flux TRS mechanism, the LC order parameter, the
Table~I multipole classification, the AHE/Chern signature, and the patch-model
channel logic. Not covered (out of scope, marked in \texttt{marker.md}): the
full weak-coupling RG flow of the patch model, the complete anharmonic Landau
free-energy minimization, and any ab initio (DFT/hybrid) prediction of LC order
in real materials.

\medskip\noindent\textbf{Agreement: 9/10.} All checked claims agree
quantitatively/qualitatively with the paper. The one point requiring care is the
Table~I dipole cancellation, which needs the symmetry-correct multipole
invariant rather than a naive scalar sum (documented). Since the paper reports no
numerical values to match digit-for-digit, agreement is judged on
correct-physics reproduction, which is complete.

\end{document}
